% Pursuing Stacks (cotes 134-2 to 134-8): archive cote 134-6, site page 21 (PDF page 22), typescript page 340.
% typescript. Draft compared: ps5.tex:2423-2477.
% AI-transcribed by Claude Opus 5.5 (claude-opus-5-5, adaptive thinking, effort high), Message Batches API, batch ps3 (msgbatch_01XdFmuVVBAcgc7EFxdtCqsj), 2026-11-12; not yet checked by eye.
\scanpage{134-6/21}
\mspage{340}

what was \unsure{heeded|needed} for recovering a 1-connected rational homotopy type
was not merely (27)  (where now L. is an object of  D.($\underline{Q}$) rather than \hand{of}
D.($\underline{Z}$) = D.(Ab)),  which reduces more or less \xout{to the} (under suitable
finiteness assumption) to knowing the rational cohomology ring,  but
an anticommutative and associative \inserted{differential graded} \uline{algebra}  over $\underline{Q}$ \over{,}{(}giving rise to (27)\inserted{by duality).}
Thus 1-connected rational homotopy types are expressible as objects of
the derived category defined in terms of such algebras,  and the \inserted[typed]{obvious} notion
of quasi-isomorphism for these.  To any space or ss set,  \deleted{he}\inserted{Sullivan} associates
a corresponding \uline{``De Rham complex''}  with rational coefficients,  in order
to get a functor from rational homotopy types to the\xout{se} derived category
obtained from those algebras  -  and (if I remember it right)  this is
an equivalence of categories,  provided one restricts to 1-connected
homotopy types\hand{,} and correspondingly to 1-connected algebras.  Probably
somebody must have explained to me by then (it was in 1976 more or less)
why not every eligible differential algebra could be recovered (up to
isomorphism\inserted{in the derived category}) by the corresponding  cohomology \xout{ring} algebra,  namely why
it was not necessarily isomorphic to the latter,  endowed with zero dif-%
ferential operator;  \xout{What} I am affraid I forgot \deleted{why}\inserted{it since} !  Also,  it was
well known by the informed people \xout{(when told} (as I was told too) that
there were obstructions against expressing the multiplicative structure
in cohomology with (say) integer coefficients,  in terms of an anticom-%
mutative differential graded\inserted{$\underline{Z}$-}algebra\over{,}{;}  so there was no hope,  I was
informed,  for defining something like a``De Rham complex with integer
coefficients'' for an arbitrary topological space.

All this \xout{sounded} was very interesting indeed  -  still I found
it hard to believe that,  while succeeding in constructing De Rham \xout{compl}
complexes with rational coefficients for arbitrary spaces,  by looking
at the algebraic De Rham complex on the enveloping affine space for
the various singular simplices of a space,  that the same could not be
achieved with integral coefficients.  Of course,  the basic Poincaré
lemma for algebraic differential forms was no longer true,  however
this reminded  me strongly of a similar difficulty met with in algebraic
geometry,  and which is overcome by working with suitable ``divided
power structures''  -  as Poincaré's lemma becomes valid when replacing
usual polynomials (as coefficients for differential forms)  by \hand{``}polyno-%
mials \unsure{``}with \unsure{devided|divided} powers''.  Then I got quite excited and involved in
\deleted{working out} a formalism of De Rham complex\deleted{es} with divided powers  for

\margin{21}
